% =====================================================================
%  Tema, cores e pacotes da apresentacao.
% =====================================================================

% --- Codificacao e idioma --------------------------------------------
\usepackage[T1]{fontenc}
\usepackage[utf8]{inputenc}
\usepackage[brazil]{babel}
\usepackage{lmodern}

% --- Elementos graficos ----------------------------------------------
\usepackage{graphicx}
\graphicspath{{figures/}}
\usepackage{booktabs}
\usepackage{amsmath}
\usepackage{amssymb}

% --- Tema ------------------------------------------------------------
% Partimos do tema mais neutro do beamer e montamos a identidade visual
% aqui, em vez de depender de um tema externo: assim a apresentacao
% compila em qualquer instalacao de TeX Live.
\usetheme{default}
\usecolortheme{default}
\useinnertheme{circles}
\setbeamertemplate{navigation symbols}{}   % remove a barra de navegacao
\setbeamertemplate{caption}[numbered]

% --- Paleta ----------------------------------------------------------
% Troque estas tres cores para adequar a apresentacao a sua instituicao.
\definecolor{lgPrimary}{RGB}{16, 58, 110}    % azul institucional
\definecolor{lgAccent}{RGB}{198, 132, 24}    % destaque
\definecolor{lgInk}{RGB}{34, 34, 34}         % texto

\setbeamercolor{normal text}{fg=lgInk}
\setbeamercolor{structure}{fg=lgPrimary}
\setbeamercolor{title}{fg=lgPrimary}
\setbeamercolor{frametitle}{fg=lgPrimary}
\setbeamercolor{alerted text}{fg=lgAccent}
\setbeamercolor{block title}{bg=lgPrimary!12, fg=lgPrimary}
\setbeamercolor{block body}{bg=lgPrimary!5}
\setbeamercolor{itemize item}{fg=lgAccent}
\setbeamercolor{itemize subitem}{fg=lgAccent}

% --- Tipografia dos titulos ------------------------------------------
\setbeamerfont{title}{size=\LARGE, series=\bfseries}
\setbeamerfont{frametitle}{size=\large, series=\bfseries}
\setbeamerfont{footline}{size=\scriptsize}

% Titulo do slide com um filete embaixo, para separar do conteudo.
\setbeamertemplate{frametitle}{%
  \vspace{0.6em}%
  \usebeamerfont{frametitle}\usebeamercolor[fg]{frametitle}\insertframetitle%
  \vspace{0.2em}%
  \hrule height 1pt \relax%
}

% --- Rodape ----------------------------------------------------------
% Autor, titulo abreviado e numero do slide, discretos.
\setbeamertemplate{footline}{%
  \hbox{%
    \begin{beamercolorbox}[wd=0.45\paperwidth, ht=2.4ex, dp=1.2ex, leftskip=1em]{author in head/foot}%
      \usebeamerfont{footline}\insertshortauthor%
    \end{beamercolorbox}%
    \begin{beamercolorbox}[wd=0.45\paperwidth, ht=2.4ex, dp=1.2ex, center]{title in head/foot}%
      \usebeamerfont{footline}\insertshorttitle%
    \end{beamercolorbox}%
    \begin{beamercolorbox}[wd=0.10\paperwidth, ht=2.4ex, dp=1.2ex, right, rightskip=1em]{date in head/foot}%
      \usebeamerfont{footline}\insertframenumber{}/\inserttotalframenumber%
    \end{beamercolorbox}%
  }%
  \vspace*{0.2em}%
}
\setbeamercolor{author in head/foot}{fg=lgInk!60}
\setbeamercolor{title in head/foot}{fg=lgInk!60}
\setbeamercolor{date in head/foot}{fg=lgInk!60}

% --- Slide de secao --------------------------------------------------
% Um slide de transicao a cada \section, para dar ritmo a apresentacao.
\AtBeginSection[]{%
  \begin{frame}[plain]
    \centering
    {\usebeamerfont{title}\usebeamercolor[fg]{title}\insertsectionhead}
  \end{frame}
}

% --- Comandos de conveniencia ----------------------------------------
% Fonte de uma figura ou tabela, em corpo reduzido.
\newcommand{\fonte}[1]{%
  \par\vspace{0.3em}{\scriptsize\color{lgInk!70}Fonte: #1}%
}
